\subsection{Functorial properties of duality}

Let G and H be locally compact commutative groups. Recall that if \(\varphi\colon\mathbf{G}\to\mathbf{H}\) is a morphism of topological groups, the dual morphism \(\widehat{\varphi}\colon\widehat{\mathrm{H}}\to\widehat{\mathrm{G}}\) is defined by \(\langle\chi,\varphi(x)\rangle=\langle\widehat{\varphi}(\chi),x\rangle\) for all \(\chi\in\widehat{\mathrm{H}}\) and \(x\in\mathrm{G}\) . This definition shows that \(\widehat{\widehat{\varphi}}=\varphi\) with the identifications of G (resp. H) and \(\widehat{\widehat{\mathrm{G}}}\) (resp. \(\widehat{\widehat{\mathrm{H}}}\) ) of theorem 2 of II, p.~220.

Let \(\theta\) be a map from \(\mathbf{G} \times \mathbf{H}\) into \(\mathbf{U}\). For every \(x \in \mathbf{G}\) (respectively, every \(y \in \mathbf{H}\)), let \(\theta_x\) (respectively, \(\theta^y\)) be the function from \(\mathbf{G}\) into \(\mathbf{U}\) defined by \(y \mapsto \theta(x, y)\) (respectively, the function from \(\mathbf{H}\) into \(\mathbf{U}\) defined by \(x \mapsto \theta(x, y)\)). Suppose that the map \(\alpha: x \mapsto \theta_x\) is an isomorphism of the topological group \(\mathbf{G}\) onto the topological group \(\widehat{\mathbf{H}}\). For every \(y \in \mathbf{H}\) and \(x \in \mathbf{G}\), one has

\[
\theta^ {y} (x) = \theta (x, y) = \langle \alpha (x), y \rangle = \langle x, \widehat {\alpha} (y) \rangle ,
\]

that is, \(\theta^{y} = \widehat{\alpha}(y)\). By th. 2 of II, p.~220, the map \(\beta: y \mapsto \theta^{y}\) is therefore an isomorphism of the topological group H onto the topological group \(\widehat{G}\). Under these conditions, we shall say that \(\theta\) puts G and H in duality, or that G and H are in duality with respect to \(\theta\). We will then identify each of the groups G and H with the dual of the other. We will call the dual measure of the Haar measure dx the Haar measure on H obtained by transport of structure from the dual measure of dx.

Lemma 7. --- Let \((\mathrm{G}_{i})_{i\in\mathrm{I}}\) and \((\mathrm{H}_{i})_{i\in\mathrm{I}}\) be finite families of locally compact topological groups. For \(i\in I\) , let \(\theta_{i}\colon G_{i}\times H_{i}\to U\) be a map that puts the groups \(G_{i}\) and \(H_{i}\) in duality. The map \(\theta\) defined by

\[
\theta ((g _ {i}), (h _ {i})) = \prod_ {i \in I} \theta_ {i} (g _ {i}, h _ {i})
\]

puts the groups \(\prod G_{i}\) and \(\prod H_{i}\) in duality.

This follows from prop. 2 of II, p.~206 and from the definition preceding it.

DEFINITION 5. --- Let G, H, and K be topological groups. Let \(f: \mathrm{H} \to \mathrm{G}\) and \(g: \mathrm{G} \to \mathrm{K}\) be morphisms of topological groups. The pair \((f, g)\) is said to be an exact sequence of topological groups if it is an exact sequence of groups (A, II, p.~10, remark 5) and if \(f\) and \(g\) are strict morphisms.

We will represent an exact sequence by the diagram

\[
\mathrm{H} \xrightarrow {f} \mathrm{G} \xrightarrow {g} \mathrm{K},
\]

and we will say that a diagram

\[
\mathrm{G} _ {1} \xrightarrow {f _ {1}} \mathrm{G} _ {2} \xrightarrow {f _ {2}} \mathrm{G} _ {3} \rightarrow \dots \rightarrow \mathrm{G} _ {n} \xrightarrow {f _ {n}} \mathrm{G} _ {n + 1}
\]

is exact if each pair \((f_i, f_{i+1})\) for \(1 \leqslant i \leqslant n-1\) is exact.

A sequence

\[
1 \to \mathrm{H} \xrightarrow {f} \mathrm{G} \xrightarrow {g} \mathrm{K} \to 1
\]

is exact if and only if f is a strict injective morphism, g is a strict surjective morphism, and the kernel of g equals the image of f. If K is separated, the image of f is a closed subgroup of G.

Examples. --- 1) Let \(f\colon H \to G\) be a strict injective morphism whose image is a distinguished subgroup. The sequence

\[
1 \to \mathrm{H} \xrightarrow {f} \mathrm{G} \xrightarrow {p} \mathrm{G} / f (\mathrm{H}) \to 1,
\]

where p is the canonical projection, is exact. In particular, if H is a closed normal subgroup of G, the sequence of topological groups

\[
1 \to \mathrm{H} \xrightarrow {j} \mathrm{G} \xrightarrow {p} \mathrm{G} / \mathrm{H} \to 1,
\]

where j is the inclusion and p the canonical projection, is exact.

\begin{enumerate}
\def\labelenumi{\arabic{enumi})}
\setcounter{enumi}{1}
\tightlist
\item
  Let \(g: G \rightarrow K\) be a strict surjective morphism. The sequence
\end{enumerate}

\[
1 \to \mathrm{Ker} (g) \stackrel {j} {\longrightarrow} \mathrm{G} \stackrel {g} {\longrightarrow} \mathrm{K} \to 1,
\]

where j is the inclusion, is exact.

Theorem 4. --- A sequence

\[
\mathrm{H} \xrightarrow {f} \mathrm{G} \xrightarrow {g} \mathrm{K}
\]

of locally compact commutative topological groups is exact if and only if the dual sequence

\[
\widehat {\mathrm{K}} \xrightarrow {\widehat {g}} \widehat {\mathrm{G}} \xrightarrow {\widehat {f}} \widehat {\mathrm{H}}
\]

is exact.

We begin by proving a few lemmas. Note that each of them is, moreover, an easy consequence of the assertion of Thm. 4.

Lemma 8. --- Let \(g \colon G \to K\) be a morphism of locally compact commutative topological groups. If the morphism \(g\) is surjective and strict, then \(\widehat{g}\) is injective and strict.

Since \(g\) is surjective, the morphism \(\widehat{g}\) is injective (lemma 2 of II, p.~205). To prove that \(\widehat{g}\) is a strict morphism, it suffices to show that for every neighborhood U of \(e\) in \(\widehat{\mathbf{K}}\), there exists a neighborhood V of \(e\) in \(\widehat{\mathbf{G}}\) such that \(\overline{\widehat{g}}^{-1} (\mathrm{V})\subset \mathrm{U}\) (lemma 2 of II, p.~200). Let U be such a neighborhood of \(e\) in \(\widehat{\mathbf{K}}\). By definition of the topology of \(\widehat{\mathbf{K}}\), there exists a compact subset X of K and a number \(\varepsilon >0\) such that U contains the set of \(\widehat{z}\in \widehat{\mathbf{K}}\) which, for every \(z\in \mathrm{X}\), satisfy \(|\langle \widehat{z},z\rangle -1| < \varepsilon\). Since \(g\) is strict and surjective, there exists, by TG, I, p.~80, prop. 10, a compact subset \(\mathrm{X}_0\) of G such that \(g(\mathrm{X}_0) = \mathrm{X}\). Let V be the neighborhood of \(e\) in \(\widehat{\mathbf{G}}\) consisting of the elements \(\chi \in \widehat{\mathbf{G}}\) such that, for all \(x\in \mathrm{X}_0\), one has \(|\langle \chi ,x\rangle -1| < \varepsilon\). We then have \(\overline{\widehat{g}}^{-1} (\mathrm{V})\subset \mathrm{U}\). This proves the assertion.

Lemma 9. --- Let \(f\colon H\to G\) be a morphism of locally compact topological groups. If the morphism \(f\) is injective and strict, then \(\widehat{f}\) is surjective and strict.

Suppose that \(f\) is injective and strict. The morphism \(\widehat{f}\) induces, by passing to the quotient, a morphism \(q \colon \widehat{\mathrm{G}} / \mathrm{Ker}(\widehat{f}) \to \widehat{\mathrm{H}}\). It remains to show that this is an isomorphism of topological groups; by duality, it suffices for this to show that its dual \(\widehat{q}\) is an isomorphism.

Let L denote the dual group of \(\widehat{\mathrm{G}}/\mathrm{Ker}(\widehat{f})\) and \(p\colon\widehat{\mathrm{G}}\to\widehat{\mathrm{G}}/\mathrm{Ker}(\widehat{f})\) the canonical projection. We shall first prove that \(\widehat{p}\) induces, by passing to subspaces, an isomorphism from L onto \(f(\mathrm{H})\) .

We have \(q \circ p = \widehat{f}\) , whence \(\widehat{p} \circ \widehat{q} = f\) . The image of \(\widehat{p}\) therefore contains \(f(\mathrm{H})\) .

Since \(f\) is strict, its image \(f(\mathrm{H})\) is a locally compact subgroup of \(\mathrm{G}\) , and is therefore closed (TG, III, p.~22, cor. 2). Let \(\mathrm{K} = \mathrm{G}/f(\mathrm{H})\) and consider the exact sequence

\[
1 \to \mathrm{H} \xrightarrow {f} \mathrm{G} \xrightarrow {g} \mathrm{K} \to 1
\]

associated with example 1. By duality, the morphism \(\widehat{f} \circ \widehat{g}\) is trivial, and hence the image of \(\widehat{g}\) is contained in \(\operatorname{Ker}(\widehat{f}) = \operatorname{Ker}(p)\) . Thus \(p \circ \widehat{g}\) is the trivial morphism and, again by duality, \(g \circ \widehat{p}\) is also trivial. It follows that the image of \(\widehat{p}\) is contained in the kernel of \(g\) , which is equal to \(f(\mathrm{H})\) . We conclude that the image of \(\widehat{p}\) is equal to \(f(\mathrm{H})\) .

Moreover, since p is a surjective and strict morphism, the dual morphism \(\widehat{p}\) is a strict injective morphism from L to G (Lemma 8). It follows that \(\widehat{p}\) induces an isomorphism of topological groups from L onto \(f(\mathrm{H})\). Since \(\widehat{p} \circ \widehat{q} = f\) and since f induces an isomorphism from H onto its image \(f(\mathrm{H})\), the morphism \(\widehat{q}\) is an isomorphism.

Lemma 10. --- Let

\[
\mathrm{H} \xrightarrow {f} \mathrm{G} \xrightarrow {g} \mathrm{K}
\]

an exact sequence of locally compact abelian groups. The kernel of \(\hat{f}\) is equal to the image of \(\hat{g}\) .

The homomorphism \(\widehat{f} \circ \widehat{g}\) is trivial by duality, so the image of \(\widehat{g}\) is contained in the kernel of \(\widehat{f}\). Conversely, let \(\chi\) be in the kernel of \(\widehat{f}\). This means that \(\operatorname{Im}(f) = \operatorname{Ker}(g)\) is contained in the kernel of \(\chi\), hence that there exists a character \(\eta\) of \(\operatorname{Im}(g)\) such that \(\eta \circ g = \chi\). Since the inclusion of \(\operatorname{Im}(g)\) in K is strict, the dual map of restriction of characters from K to \(\operatorname{Im}(g)\) is surjective (Lemma 9). There therefore exists a character \(\beta\) of K such that \(\eta\) is the restriction of \(\beta\), and it follows that \(\chi = \beta \circ g = \widehat{g}(\beta)\). From this we conclude that the kernel of \(\widehat{f}\) is contained in the image of \(\widehat{g}\).

Now let us prove Theorem 4. By duality, it suffices to show that the sequence \(\widehat{\mathrm{K}}\xrightarrow{\widehat{g}}\widehat{\mathrm{G}}\xrightarrow{\widehat{f}}\widehat{\mathrm{H}}\) is exact when the sequence \(\mathrm{H}\xrightarrow{f}\mathrm{G}\xrightarrow{g}\mathrm{K}\) is exact. Now, by Lemmas 8 and 9, the morphisms \(\widehat{f}\) and \(\widehat{g}\) are strict and, by Lemma 10, the kernel of \(\widehat{f}\) is equal to the image of \(\widehat{g}\).

COROLLARY 1. --- Let

\[
1 \to \mathrm{H} \xrightarrow {f} \mathrm{G} \xrightarrow {g} \mathrm{K} \to 1
\]

an exact sequence of locally compact commutative topological groups. The morphism \(\widehat{g}\) induces an isomorphism between \(\widehat{\mathrm{K}}\) and \(f(\mathrm{H})^{\perp}\) , and \(\widehat{f}\) induces, by passing to the quotient, an isomorphism between \(\widehat{\mathrm{G}} / f(\mathrm{H})^{\perp}\) and \(\widehat{\mathrm{H}}\) .

By Theorem 4, the sequence

\[
1 \to \widehat {\mathrm{K}} \xrightarrow {\widehat {g}} \widehat {\mathrm{G}} \xrightarrow {\widehat {f}} \widehat {\mathrm{H}} \to 1\tag{32}
\]

is exact. The morphism \(\widehat{g}\) therefore induces an isomorphism from \(\widehat{\mathbf{K}}\) onto \(\mathrm{Ker}(\widehat{f}) = f(\mathrm{H})^{\perp}\) (lemma 2 of II, p.~205), and \(\widehat{f}\) induces, by passing to the quotient, an isomorphism from \(\widehat{\mathrm{G}} / \mathrm{Ker}(\widehat{f}) = \widehat{\mathrm{G}} / f(\mathrm{H})^{\perp}\) onto \(\widehat{\mathrm{H}}\) (loc. cit.).

COROLLARY 2. --- Let \(f\colon G\to H\) be a morphism of locally compact commutative topological groups. The morphism \(f\) is strict if and only if \(\widehat{f}\) is strict.

By the canonical decomposition (E, II, p.~44) of a strict morphism, this follows from Lemmas 8 and 9.

COROLLARY 3. --- Let H be a subgroup of G. Then (H \(^{\perp}\) ) \(^{\perp}\) = \(\overline{H}\) .

Since \(H^{\perp} = \overline{H}^{\perp}\) , we may assume that H is closed. Let \(f: H \to G\) be the canonical injection and \(p: G \to G/H\) the canonical projection. We have \(H^{\perp} = \operatorname{Ker}(\widehat{f})\) (Lemma 10). Let k be the canonical injection of \(H^{\perp}\) into \(\widehat{G}\). By Theorem 4, the morphism \(\widehat{p}\) induces an isomorphism \(\widehat{p}_{H}: \widehat{G/H} \to H^{\perp}\) of topological groups and we have \(k \circ \widehat{p}_{H} = \widehat{p}\). Consequently (Corollary 1), we obtain

\[
(\mathrm{H} ^ {\perp}) ^ {\perp} = \mathrm{Ker} (\widehat {k}) = \mathrm{Ker} (\widehat {\widehat {p}} _ {\mathrm{H}} \circ \widehat {k}) = \mathrm{Ker} (p) = \mathrm{H}.
\]

COROLLARY 4. --- Let I be a set and let \((\mathrm{H}_{i})_{i\in\mathrm{I}}\) be a family of closed subgroups of G. The orthogonal of the closed subgroup generated by the \(H_{i}\) is \(\bigcap_{i\in I}H_{i}^{\perp}\) . The orthogonal of \(\bigcap_{i}H_{i}\) is the closed subgroup generated by the subgroups \(H_{i}^{\perp}\) .

The first assertion follows from the definition of the orthogonal. Applying this result and Cor. 3 to the family of closed subgroups \((\mathrm{H}_{i}^{\perp})_{i\in\mathrm{I}}\) of \(\widehat{G}\) we see that \(\bigcap_{i}H_{i}\) is the orthogonal of the closed subgroup generated by the subgroups \(H_{i}^{\perp}\) and the second assertion is then obtained by duality.

COROLLARY 5. --- Let \(\varphi\) : G → H be a morphism of locally compact commutative groups. Then the subgroup \(\overline{\text{Im}(\varphi)}\) of H and the subgroup \(\text{Ker}(\widehat{\varphi})\) of \(\widehat{\text{H}}\) are orthogonal to each other. In particular, for \(\widehat{\varphi}\) to be injective, it is necessary and sufficient that the image of \(\varphi\) be dense in H.

We have \(\operatorname{Ker}(\widehat{\varphi}) = \varphi(\mathrm{G})^{\perp}\) (lemma 2 of II, p.~205), whence the result by cor. 3.

COROLLARY 6. --- Let \(k \in \mathbf{Z}\). Then the kernel of the homomorphism \(x \mapsto x^k\) from G to G and the closure of the image of the morphism \(\chi \mapsto \chi^k\) from \(\widehat{\mathrm{G}}\) to \(\widehat{\mathrm{G}}\) are orthogonal to each other.

This follows from the preceding corollary, since the morphisms \(x \mapsto x^{k}\) from G to G and \(\chi \mapsto \chi^{k}\) from \(\widehat{G}\) to \(\widehat{G}\) are dual to each other.

Recall (A, X, p.~17) that a commutative group A is divisible if, for every nonzero \(n \in \mathbf{Z}\), the map \(a \mapsto a^n\) from A to A is surjective.

COROLLARY 7. --- Let G be a locally compact abelian group.

\begin{enumerate}
\def\labelenumi{\alph{enumi})}
\item
  If G is divisible, then the dual group \(\widehat{\mathrm{G}}\) is torsion-free;
\item
  If the dual group \(\widehat{\mathbf{G}}\) is torsion-free, and if \(k\in\mathbf{Z}\) is nonzero, then the image of the homomorphism \(x\mapsto x^{k}\) from G to G is dense in G;
\item
  Suppose G is discrete or compact. For G to be divisible it is necessary and sufficient that \(\widehat{G}\) be torsion-free.
\end{enumerate}

The assertions \(a\) ) and \(b\) ) follow from cor. 6. If G is discrete or compact, the image of the homomorphism \(x \mapsto x^k\) from G to G is closed, and \(c\) ) follows from \(a\) ) and \(b\) ).

Remark. --- There exist locally compact abelian groups G that are not divisible and such that \(\widehat{G}\) is torsion-free (exercise 63 of II, p.~299).

